\chapter{Advanced Calculus - 1}
\section{}

% Set the course for this chapter
\Course{Advanced Calculus - 1}

\section{Hafta 1: Vektörler ve Temel Kavramlar}
\Week{1}
\begin{definition}
    \Topic{Function of Several Variables}
A function of two variables is a rule that assigns a real number
We sometimes write: $f: D \subset \mathbb{R}^2 \to \mathbb{R}$; for instance, $f(x,y) = x^2 + y^2$ or $g(x,y) = x^2 - e^y$.

    Likewise, a function of three variables is a rule that assigns a real number to each ordered triple
     $(x,y,z)$ in its domain $D \subset \mathbb{R}^3$. We sometimes write: $f: D \subset \mathbb{R}^3 \to \mathbb{R}$;
      for instance, $f(x,y,z) = x^2 + y^2 + z^2$ or $g(x,y,z) = x^2 - e^{y+z}$.

      $f : D \subset \mathbb{R}^n \to \mathbb{R}$ is a function of $n$ variables if it assigns a real number to each
       ordered $n$-tuple $(x_1, x_2, \ldots, x_n)$ in its domain $D \subset \mathbb{R}^n$.
\end{definition}



\begin{exampleauto}
Find the domain of the functions.

\begin{example}
    \topic{Example 1.1}
    $f(x,y) = \ln{4 - x^2 - 4y^2}$
\end{example}
    
\begin{example}
    \topic{Example 1.2}
    $g(x,y) = \sqrt{x^2 + y^2 -1} + \sqrt{4 - x^2 - y^2}$
\end{example}
 
\begin{example}
     \topic{Example 1.3}
    $f(x,y) = \frac{1}{\sqrt{4x^2 - y^2}}$
\end{example}

\end{exampleauto}




\begin{definition}
    \Topic{Neighborhoods in $\mathbb{R}^n$}
A neighborhood of a point $\mathbf{a} \in \mathbb{R}^n$ is a set that contains all points $\mathbf{x}$ that are "close" to $\mathbf{a}$. More formally, we can define a neighborhood of $\mathbf{a}$ as follows:

Let $r > 0$. The open ball of radius $r$ centered at $\mathbf{a}$ is defined as:
\[
B_r(\mathbf{a}) = \{ \mathbf{x} \in \mathbb{R}^n : \|\mathbf{x} - \mathbf{a}\| < r \}
\]
A neighborhood of $\mathbf{a}$ is any set that contains an open ball centered at $\mathbf{a}$.
\end{definition}

\begin{definition}
    \Topic{The interior of a set}
The interior of a set $A \subset \mathbb{R}^n$, denoted $\operatorname{int}(A)$, is the largest open set contained in $A$. In other words, a point $\mathbf{x} \in A$ is in the interior of $A$ if there exists a neighborhood of $\mathbf{x}$ that is entirely contained in $A$.
\end{definition}

\begin{definition}
    \Topic{The boundary of a set}
The boundary of a set $A \subset \mathbb{R}^n$, denoted $\partial A$, is the set of points that can be approached both from inside $A$ and from outside $A$. Formally, a point $\mathbf{x} \in \mathbb{R}^n$ is in the boundary of $A$ if every neighborhood of $\mathbf{x}$ contains at least one point in $A$ and at least one point not in $A$.
\end{definition}